% !TeX root = ../papers/nc-p3.tex
\section{十五次元成分内のWalsh三次元族とその外の周期層}
\label{sec:walsh-full-component}
\begin{theorem}\label{thm:walsh-full-component}
Walsh三次元族の全体は、添字保持同型を除いた次元が15の
定理\ref{thm:full-fifteen-component}の既約成分に属する。
一般のWalsh点で隣接階数スキームは滑らかである。指定なしの成分内で、
この点を通る周期が2を割る軌跡の被約局所枝は四つあり、次元は $9,5,5,5$ である。
各枝の幾何学的一般点は標準Koszul左右AS正則で最小周期2を持ち、
Walsh族の外にある。残る一般の変形は最小周期4である。
\end{theorem}
\begin{proof}[二つの表示族の交点の滑らかさ]
$a=1,b=3,r=4$, $l^2=-10$, $p=5$, $(m')^2=-1$, $q^2=17$, $n^2=-9$ とし、
53の上で $(l,m',q,n)=(19,23,21,16)$ となる根を選ぶ。
階数15の復元式では $(s,t,N')=(25,35,28)$ に還元される。通常な正の固定曲線の第一表示で
倍率 $(z,\lambda)=(15,38)$、枠 $(a',b',e',f')=(29,1,26,17)$ となる。
定理\ref{thm:klein-elliptic-completion}の二次式 $P,Q$ の先頭係数は剰余で
$4,36$、選んだ根での導関数は $7,-7$、二枠の行列式は $31,51$ である。
従ってこの標数零の窓は同定理の既約な二重被覆の対根表示に実際に含まれる。
分母が単元で、根を選ぶモニックな二重被覆の一般点をこの点へ特殊化できる。
従って既知の十五次元成分上のWalsh点である。

この点の第一辺は53の上で二十個の横断交点を持つ。四座標点と、
$(1,3,10,2)$ に四ラベルの平行移動・文字倍を施した十六点を使う。
各点で次の点の四座標は一次元核で決まり、六双一次式のJacobi階数は6である。
Segre多様体の次数は20なので、この二十個の孤立単純交点が全交叉を尽くす。
実際、線型切断を順に取ったときの全既約成分の次数の和は増えず、
正次元成分があればその次数も加算される。各点での二つの部分評価は階数3で、
得た核直線を用いる定理\ref{prop:twenty-conormal}の二十行は独立である。
\path{computations/verify_walsh_full_component.py} はこの計算と以下の接計算を再現する。
有限エタールな二十点と単元小行列によってこれらの条件は標数零に持ち上がる。

二切断の接方程式は、それぞれ摂動の両側に十次元の左右核を掛けて得る。
位数16の $G$ の文字を $(h,c)$ と書く。$h$ は語の総色、$c$ は同時移動に対する文字で、
$g=(x,z)$ は符号 $(-1)^{h\cdot z+c\cdot x}$ で作用する。
53での接空間の文字次元は $h=0$ の四文字で各4、$h\ne0$ の十二文字で各5である。
従って接方程式の階数は180以上で、二十恒等式により180以下でもある。
命題\ref{prop:twenty-conormal}で非射影化したスキームは滑らかな76次元となる。
文字次元の上界の合計が76なので個々の値も標数零で一致する。
滑らかな点では成分が一意であるから、この点を通る既約なWalsh族の全体が同じ成分に入る。
その一般点でも滑らかさと文字次元は保たれる。

アフィン軌道の接空間では自明文字は倍率の一つ、各非自明文字は四位置の行列から一つずつである。
実際各位置のPauli共役表現は行列空間で十六文字を一度ずつ持ち、
一般Walsh点の射影安定化群は有限である。
従って横断する十五次元空間は、自明表現三つと
$h\ne0,c\in(\mathbb Z/2)^2$ の十二の文字表現に分かれる。
\end{proof}

\begin{proof}[二位置移動の固定部分]
一般のWalsh点は幾何学的に安定で射影安定化群は $G$ である。
実際、正規直交基底に関するノルム最小性と安定化群の有限性は
定理\ref{thm:walsh-presentation-space}で確認した。二位置移動に四つの
対角変換 $J(i)=(-1)^{i_0 i_1}$ を合成して、表示点を固定する対合 $\tau$ を得る。
$\tau$ は四基底変更群を正規化する。安定軌道のエタール切片定理を、
二位置移動も加えた有限拡大の群に適用する
（D.~Luna, \emph{Slices étales}, Mémoires de la S.M.F. 33, 81--105, 1973）。
滑らかな横断切片の形式座標は有限安定化群で平均して線型化できる。
元の指定なし商の局所模型は上の十五次元空間の $G$ 商で、
その二位置移動には $G\rtimes\langle\tau\rangle$ の作用が誘導される。

$\tau$ の共役は $G$ の $(x,z)$ を $(x,z+x^\flat)$ に移す。
ここで $x^\flat=(x_1,x_0)$ であり、これは $J(i+x)/J(i)$ の線型部分である。
従って $h\ne0$ では $\tau$ が $(h,c)$ と $(h,c+h^\flat)$ の一次元方向を交換する。
各対の基底を選べば係数1の交換にできる。
三つの自明方向はWalsh族の接方向で、族の表示全体を $\tau$ が固定するので作用は恒等である。

有限商での移動不動点は、切片上の各 $g\tau$ の固定線型部分空間の像の合併である。
$g=(x,z)$ とし、各対の二座標を $A,B$ と書く。
$x\cdot h^\flat=0$ なら $B=(-1)^{h\cdot z+c\cdot x}A$、
そうでなければ両方零となる。$x=0$ なら六対の全部が残り、
各 $h$ で二対の符号が同じである。これら四つの九次元空間は $G$ の共役で一つの枝になる。
$x\ne0$ なら残るのは $h=x$ の二対だけであり、両対の符号は互いに逆になる。
二符号を同時に反転する二つの五次元空間は $G$ で同一視される。
三つの $x\ne0$ は異なる文字の支持を持つので互いに同一視されない。
また逆の二符号を持つ五次元空間は、同じ二符号を持つ九次元空間のいずれにも含まれない。
これで被約局所枝が尽きる。有限群の線型空間の軌道から来る各枝は
形式的にも既約である（その有限正規化の該当点は一軌道である）。

安定商の分離性により一位置移動の不動点は閉で、一般Walsh点はそこに属さない。
各局所枝の幾何学的一般点の全次数の正則性は、正則な特殊点から
有限Frobenius双対の内部次数別bar複体を持ち上げて得る。
各枝はWalsh三次元より大きく、成分の一般点には二位置移動も残らない。
\end{proof}
